
\documentclass[12pt]{article}
\usepackage{amsmath} % AMS Math Package
\usepackage{amsthm} % Theorem Formatting
\usepackage{amssymb}    % Math symbols such as \mathbb
\usepackage{graphicx} % Allows for eps images
\usepackage[dvips,letterpaper,margin=1in,bottom=0.7in]{geometry}
\usepackage{tensor}
\newtheorem{p}{Problem}[section]
\usepackage{amsmath}

\begin{document}

{\noindent\Huge\bf  \\[0.5\baselineskip] {\fontfamily{cmr}\selectfont  Problem Set II}         }\\[2\baselineskip] % Title
{ {\bf \fontfamily{cmr}\selectfont Political Science 43401}\\ {\textit{\fontfamily{cmr}\selectfont     October 1, 2020}}}\hfill   {\large \textsc{Robert Gulotty}} % Author name
\\[1.4\baselineskip] 





\begin{enumerate}

\item[5.1.3] 
\item[5.2.1]
\item[6.3.2] (p. 114)
\item[ 6.3.3] (p. 114)
\item[ 7.1.1] (b) (p. 121)
\item[7.1.2] (b) (p. 121)
\item[7.2.6] (p. 125)
\item[ 7.3.2] (a) (p. 130)
\item[7.4.1]
\item[7.4.3]
\item[31.1.3]
\item[31.3.3]
\item[31.4.2 ]
\item[31.4.5 ] 
\item[R bonus] Using R, replicate Figure 6.5
\end{enumerate}

\end{document}